% titel: Blatt 9, A11: L von Ag2CrO4 und Löslichkeit von AgBr
% Teil II: der durchgerechnete Weg.

\subsection*{Gegeben und gesucht}
a) $c(Ag_2CrO_4)=7.8\cdot10^{-5}\ mol\,L^{-1}$ (25\,°C). Gesucht: $L$.
b) $L(AgBr)=4.0\cdot10^{-13}\ mol^2L^{-2}$; $c(NaBr)=0.10\ M$. Gesucht:
molare Löslichkeit $s$ von AgBr.

\subsection*{a) Ag$_2$CrO$_4$}
Aus $Ag_2CrO_4$ entstehen $2\,Ag^+$ und ein $CrO_4^{2-}$, also
$[Ag^+]=2s$, $[CrO_4^{2-}]=s$ mit $s=7.8\cdot10^{-5}$.
\[
  L = [Ag^+]^2[CrO_4^{2-}] = (2s)^2\,s = 4s^3
    = 4\,(7.8\cdot10^{-5})^3 .
\]
\[
  (7.8\cdot10^{-5})^3 \approx 4.74\cdot10^{-13},
\]
\[
  L \approx 4\cdot4.74\cdot10^{-13} \approx 1.9\cdot10^{-12}\ mol^3L^{-3}.
\]

\subsection*{b) AgBr mit gemeinsamen Ion}
AgBr dissoziiert zu Ag$^+$ und Br$^-$ (1:1). In 0.10 M NaBr ist
$[Br^-]\approx0.10$ (Fremdion, praktisch konstant):
\[
  L = [Ag^+][Br^-]\quad\Rightarrow\quad
  s = [Ag^+] = \frac{L}{[Br^-]} = \frac{4.0\cdot10^{-13}}{0.10}.
\]
\[
  s \approx 4.0\cdot10^{-12}\ mol\,L^{-1}.
\]

\begin{ergebnis}
a) $L(Ag_2CrO_4)\approx1.9\cdot10^{-12}$
$\mathrm{mol^3L^{-3}}$.
b) Löslichkeit von AgBr in 0.10 M NaBr: $s\approx4.0\cdot10^{-12}$
mol$\,L^{-1}$.
\end{ergebnis}

\begin{kontrolle}
a) $4\cdot(7.8e{-}5)^3 = 1.9e{-}12$, Probe ok. b) Ohne NaBr wäre
$s=\sqrt{L}\approx2.0e{-}7$; mit 0.10 M Fremdion $4e{-}12\ll$ --- der
gemeinsame Ioneneffekt unterdrückt die Löslichkeit massiv, wie erwartet.
Konsistent.
\end{kontrolle}

\begin{merksatz}
$L$ an die Stöchiometrie koppeln ($4s^3$ bei 2:1). Gemeinsames Ion setzt
$s = L[\text{Anion}]$ durch, die Löslichkeit sinkt drastisch.
\end{merksatz}
